\documentclass[11pt]{article}
\usepackage{preamble}

\newcommand{\IconGraph}{[GRAPH]}
\newcommand{\IconNote}{[NOTE]}
\newcommand{\lessonrow}[2]{\textbf{#1} & #2 \\ \hline}
\newcommand{\rulescardcontent}{%
LOG \(\Leftrightarrow\) EXP: \ \ \(\log_b x = y\) \ \(\Leftrightarrow\)\ \ \(b^y = x\)\\
COMMON LOG: \ \(\log x = \log_{10} x\)\\
NATURAL LOG: \(\ln x = \log_e x\)\\
KEY: \ \ \ \ \ \ \ \ \(\ln(e^t) = t\) \ \ \ \ (this solves continuous growth!)\\
UNDEFINED: \(\log_b 0\) and \(\log_b(\text{negative})\) --- no real answer\\
IDENTITY: \ \ \(\log_b(b^k) = k\)%
}
\newcommand{\qfifteengraph}{%
\begin{center}
\begin{tikzpicture}
\begin{axis}[
  width=0.88\linewidth,
  height=2.75in,
  xmin=-1,
  xmax=4.5,
  ymin=0,
  ymax=85,
  xtick={-1,0,1,2,3,4},
  ytick={0,10,20,30,40,50,60,70,80},
  xlabel={$x$},
  ylabel={$y$},
  grid=major,
  minor tick num=1,
  tick label style={font=\small},
  label style={font=\small},
  every axis plot/.append style={line width=1pt}
]
\addplot[tealheader, smooth, domain=-1:4.0, samples=250] {3^x};
\addplot[only marks, mark=*, mark size=2pt, color=dayblue] coordinates {
  (-1,0.3333333333)
  (0,1)
  (1,3)
  (2,9)
  (3,27)
  (4,81)
};
\node[anchor=south west, font=\small] at (axis cs:3.18,33) {$y=3^x$};
\end{axis}
\end{tikzpicture}
\end{center}
}
\newcommand{\richterplaceholder}{%
\begin{center}
\fbox{%
\parbox{0.9\linewidth}{%
\textbf{[IMAGE:} Cracked wall photo with Richter scale overlay numbered 1 through 10 and labels ``Not felt'' at 1 and ``Extreme'' near 9--10; callout text: ``At a richter magnitude of 4 and above, the walls in your house may start to crack.''\textbf{]}%
}%
}
\end{center}
}

\begin{document}

\packetheading{Lesson 6-3 \textperiodcentered\ Quick \textperiodcentered\ Student Packet}{Logarithms + DOK-3 Spine: q15 (exp\(\leftrightarrow\)log inverse via graph)}

\sectionbanner{OBJECTIVES}
{\small
\renewcommand{\arraystretch}{1.25}
\noindent\begin{tabular}{|p{1.8in}|p{4.8in}|}
\hline
\lessonrow{Math Objective}{Evaluate logarithms using the definition \(\log_b x = y \Leftrightarrow b^y = x\), distinguish common (\(\log\)) vs natural (\(\ln\)) logs, and use the graph of \(y=3^x\) to estimate logarithmic values through the exp\(\leftrightarrow\)log inverse relationship.}
\lessonrow{Language Objective}{Explain the inverse relationship between exponential and logarithmic equations using ``undo'' language.}
\lessonrow{Essential Understanding}{A logarithm answers the question: what exponent? -- logs undo exponentials. If I know how to graph \(y = 3^x\), I can use that same picture to estimate \(\log_3 50\) because logarithms and exponentials are inverses.}
\end{tabular}
}

\sectionbanner{TOPIC GOALS / ESSENTIAL QUESTION / MATERIALS}
{\small
\renewcommand{\arraystretch}{1.25}
\noindent\begin{tabular}{|p{1.8in}|p{4.8in}|}
\hline
\lessonrow{Topic Goals}{Fluency evaluating logarithms; use the exp\(\leftrightarrow\)log inverse relationship graphically (Topic 6 Form B prep: q15 covers Q7, Q10, Q11; q17 covers Q13).}
\lessonrow{Essential Question}{If I know how to graph \(y = 3^x\), how can I use that same picture to estimate \(\log_3 50\) --- and why does this work?}
\lessonrow{Materials}{Student packet with printed \(y=3^x\) graph (\(x \in [-1, 4.5]\), \(y \in [0, 85]\), gridlines). Teacher copy adds dashed \(y=50\) line and read-off annotation. [GRAPH PLACEHOLDER --- needs\_cleanup=YES, see Visuals Checklist].}
\end{tabular}
}

\sectionbanner{LAUNCH --- Evaluate Logarithms (teacher think-aloud Ex 3)}

\begin{callout}{\IconBook\ LOG \(\leftrightarrow\) EXP RULES (keep printed on your page)}{calloutgreen}
\rulescardcontent
\end{callout}

\sectionbanner{EXPLORE --- Think-Pair-Share (25 min)}
\textit{Work with a partner. Start with \IconStar\ Practice \#15 (DOK-3 spine --- 10 min TPS + 5 min pair presentations). Then move to \#52, \#53, \#54, and \#17.}

\begin{callout}{\IconGraph\ GRAPH FOR q15 --- \(y = 3^x\) [GRAPH PLACEHOLDER]}{calloutyellow}
[GRAPH PLACEHOLDER --- print the \(y=3^x\) graph here: \(x \in [-1, 4.5]\), \(y \in [0, 85]\), with gridlines. Source: a2\_6-3\_SE.pdf near end of practice problems. See Visuals Checklist --- needs\_cleanup=YES.]

\qfifteengraph
\end{callout}

\textbf{Bank-item labels (in order):}
\begin{enumerate}[leftmargin=1.6em,itemsep=2pt,topsep=2pt]
  \item Practice \#15 \IconStar\ DOK-3 SPINE --- Higher Order Thinking: Use \(y=3^x\) graph to estimate \(\log_3 50\)
  \item Practice \#52 --- DOK-2 Application
  \item Practice \#53 --- DOK-2 Application [45-min variant: skip]
  \item Practice \#54 --- DOK-2 Application [45-min variant: skip]
  \item Practice \#17 --- DOK-3 Proof-style: critique \(\log_3\left(\frac{1}{27}\right) = -3\) (supports Form B Q13)
\end{enumerate}

\bankitem{Practice \#15 \IconStar\ DOK-3 SPINE --- Higher Order Thinking: Use \(y=3^x\) graph to estimate \(\log_3 50\)}{%
Higher Order Thinking Use the graph of \(y=3^x\) to estimate the value of \(\log_3 50\). Explain your reasoning.

\textit{[GRAPH / TIKZ figure]}
}{}

\bankitem{Practice \#52 --- DOK-2 Application}{%
How long does it take for \$250 to grow to \$600 at 4\% annual percentage rate compounded continuously? Round to the nearest year.
}{}

\bankitem{Practice \#53 --- DOK-2 Application [45-min variant: skip]}{%
Model with Mathematics Michael invests \$1{,}000 in an account that earns a 4.75\% annual percentage rate compounded continuously. Peter invests \$1{,}200 in an account that earns a 4.25\% annual percentage rate compounded continuously. Which person's account will grow to \$1{,}800 first?
}{}

\bankitem{Practice \#54 --- DOK-2 Application [45-min variant: skip]}{%
Reason The Richter magnitude of an earthquake is \(R = 0.67\log(0.37E) + 1.46\), where \(E\) is the energy (in kilowatt-hours) released by the earthquake.

\richterplaceholder

\textbf{a.} What is the magnitude of an earthquake that releases 11{,}800{,}000{,}000 kilowatt-hours of energy? Round to the nearest tenth.

\textbf{b.} How many kilowatt-hours of energy would earthquake have to release in order to be an 8.2 on the Richter scale? Round to the nearest whole number.

\textbf{c.} What number of kilowatt-hours of energy would an earthquake have to release in order for walls to crack? Round to the nearest whole number.
}{}

\bankitem{Practice \#17 --- DOK-3 Proof-style: critique \(\log_3\left(\frac{1}{27}\right) = -3\) (supports Form B Q13)}{%
Use Structure A student says that \(\log_3\left(\frac{1}{27}\right)\) simplifies to -3. Is the student correct? Explain.
}{}

\begin{callout}{\IconNote\ LOG \(\leftrightarrow\) EXP RULES --- Reference Card (print this!)}{calloutyellow}
\rulescardcontent
\end{callout}

\sectionbanner{SHARE / SUMMARY (5 min)}

\begin{sentenceframebox}
\textbf{Sentence frames}

Pair share (q15): On the graph of \(y = 3^x\), I found \(\log_3 50\) by looking for \(y =\) \blank[0.5in] and reading the x-value. Since \(3^3 =\) \blank[0.45in] and \(3^4 =\) \blank[0.45in], the answer is between \blank[0.45in] and \blank[0.45in], closer to \blank[0.65in]. This works because \(\log_3 50 = x\) means \(3^x =\) \blank[0.45in], so the exponential and logarithm are \blank[1in].

Whole class: one pair explains why reading the graph works. Class signals thumbs up/sideways.
\end{sentenceframebox}

\begin{callout}{\IconSearch\ BRIDGE TO LESSON 6-4}{calloutpeach}
Lesson 6-4 shifts to logarithmic FUNCTIONS --- graphing \(y = \log_b x\), identifying domain, range, and asymptotes. Today's log \(\leftrightarrow\) exp definition is the foundation.
\end{callout}

\sectionbanner{EXIT}

\begin{summaryexitbox}{EXIT TICKET --- Summary Recap}
\(\log_b x = y\) means \blank[1.5in].

\(\log_2 32 =\) \blank[0.5in] because \(2^{?} = 32\).

To solve \(e^t = 6.125\), take the \blank[0.9in] of both sides giving \(t =\) \blank[0.8in].

One biggest thing I learned today: \blank[2.7in]
\end{summaryexitbox}

\clearpage

\sectionbanner{OPTIONAL REINFORCEMENT (not collected)}
\textit{If you finish Explore early. \#18 asks you to evaluate whether \(\ln 1000 = 3\) --- use the log \(\leftrightarrow\) exp definition and the value of \(e\) to verify or refute.}

\bankitem{Optional \#1 \ \ (Savvas Practice \#18 --- \(\ln 1000 \neq 3\)?)}{%
Use Structure Explain why the expression \(\ln 1{,}000\) is not equal to 3.
}{}

\end{document}
